\documentclass[10pt,a4paper]{amsart}
\usepackage[margin=19mm]{geometry}
\usepackage{amsmath,amssymb,mathtools}
\usepackage[scr=rsfs,cal=euler]{mathalfa}
\usepackage{xfrac}
\usepackage[unicode,hidelinks,bookmarks=false]{hyperref}
\newcommand{\ie}{i.e.\ }
\newcommand{\cf}{cf.\ }
\newcommand{\resp}{respectively\ }
\newcommand{\Subsection}{\subsection}
\providecommand{\nobreakdash}{\nobreak\hbox{-}}
\setlength{\emergencystretch}{2em}
\multlinegap0pt
\usepackage{bm}
\usepackage{bbm}
\usepackage{dsfont}
\usepackage{xspace}
\usepackage{cancel}
\usepackage{mathtools}
\usepackage{amsthm}
\makeatletter
\@ifundefined{theorem}{\newtheorem{theorem}{Theorem}}{}
\@ifundefined{lemma}{\newtheorem{lemma}[theorem]{Lemma}}{}
\makeatother
\title[Hessian Geometry of Latent Space in Generative Models]{Hessian Geometry of Latent Space in Generative Models}
\author{Alexander Lobashev, Dmitry Guskov, Maria Larchenko, Mikhail Tamm}
\date{Source: arXiv:2506.10632v1 (CC BY 4.0)}
\begin{document}
\maketitle

\noindent\emph{Source and licence.} Hessian Geometry of Latent Space in Generative Models. Version arXiv:2506.10632v1 (CC BY 4.0), distributed under the Creative Commons Attribution 4.0 International licence, as recorded in the arXiv licence metadata for this exact version and shown on the abstract page https://arxiv.org/abs/2506.10632v1. Full rights notice: licenses/arxiv-2506.10632-CC-BY-4.0.txt.\par\smallskip

\noindent\textit{Editorial scope.} Counted unit: the Theorem whose source label is \texttt{None} in the article (source0.tex lines 618--638, source label \texttt{None}), delivered together with the article's own complete original proof of it (source0.tex lines 640--730, one \texttt{proof} environment of 91 lines). No mathematics is written, repaired or re-derived: statement and proof are the article's own text line for line. Required context retained uncounted: Lemma:Integral (lines 734--744). Every definition, display and label of those lines is retained unchanged. Omitted from this package and not redistributed: 1--617, 639--639, 731--733, 745--1305. Nothing inside a retained excerpt is omitted or altered: every retained body is delivered exactly as the article's lines are written, so no omission receipt of any kind accompanies this package. Citations: the retained excerpts carry no citation command at all, so no bibliography entry of the article is transcribed and no reference is redistributed. Reference adaptation: none was needed and none was made; every reference inside the delivered unit resolves inside it, so no unresolved reference or citation is produced. Printed slips kept literal: no printed word, symbol, hypothesis or number of the counted statement or of its proof was altered. Layout and class: the article's own class is \texttt{article} with packages including microtype, graphicx, subfigure, booktabs, amsmath, amsfonts, tabularray, hyperref, icml2025, amssymb, mathtools, amsthm, cleveref, todonotes; the wrapper is the lane's pinned \texttt{amsart} editorial wrapper at 10pt on A4, which restates the theorem-like environments the retained text needs and adds the article's own macro definitions that the retained text needs (none), with extra wrapper packages (bm, bbm, dsfont, xspace, cancel, mathtools). The delivered numbering is the unit's own. Assets: no retained span contains a figure environment, a graphics-inclusion command, a PDF-page inclusion, a table or a TikZ picture; no image file is copied into this package, the package declares no asset dependency and no image file is redistributed with it. Licence: version arXiv:2506.10632v1 of this article is distributed under the Creative Commons Attribution 4.0 International licence, as recorded in the arXiv licence metadata for this exact version and shown on the abstract page https://arxiv.org/abs/2506.10632v1. A full rights notice is delivered at licenses/arxiv-2506.10632-CC-BY-4.0.txt. Characters: every retained excerpt is pure ASCII, so the delivered engine, XeLaTeX with its default Latin Modern text font, reports no missing character.
\begin{lemma}\label{Lemma:Integral}
Let $X_i: \Omega \to \mathbb{R}^n$ be measurable functions, $s \in S$, $\omega \in \Omega$, $N \in \mathbb{N}$, and assume $\mathbb{E}[X_i] = \mu_x < \infty$ (first moment exists). Define:
\begin{equation}
\phi(s, N, \omega) = \psi\left(s, \frac{1}{N} \sum_{i=1}^N X_i(\omega)\right),
\end{equation}
where $\psi(x,s)$ is a continuous function in both variables and uniformly continuous in $x$ for all $s$.
Then:
\begin{equation}
I \overset{\text{def}}{=} \lim_{N \to \infty} \frac{1}{N} \log \left( \int_S e^{N \psi(s, \frac{1}{N} \sum_{i=1}^N X_i)} ds \right) \overset{\text{a.s.}}{=} \max_{s \in S} \psi(s, \mu_x) \overset{\text{def}}{=} M.
\end{equation}
\end{lemma}

\begin{theorem}
Let $ \mathcal{X} $ be a space of data samples $ x \in \mathcal{X} $, and $ S \subset \mathbb{R}^n $ be a compact domain with the continuous prior distribution $p(t)$ supported on $S$.
Suppose the conditional distribution of data samples given parameter $t$ is an exponential family
\begin{equation}
p(x|t) = e^{\langle t, f(x) \rangle - \log Z(t)},
\end{equation}
where 
\begin{equation}
Z(t) = \int_{\mathcal{X}} e^{\langle t, f(x) \rangle} dx
\end{equation}
converges for all $ t \in S $. Let $ x_1, \ldots, x_N \sim p(x|t') $. Then, as $N\to\infty$ the posterior distribution satisfies:
\begin{equation}
 \lim_{N \to \infty} \left(p(t|x_1, \ldots, x_N)\right)^{1/N} \overset{\text{a.s.}}{=} e^{-D_\text{log Z(t)}(t, t')} = e^{-D_\text{KL}(p(x|t') \| p(x|t))}.
\end{equation}
where $D_\text{log Z(t)}(t, t')$ is the Bregman divergence between exponential family distributions
\begin{equation}
    \begin{split}
        &D_\text{log Z(t)}(t, t') = \log Z(t) - \log Z(t') - \langle \nabla_{t'} \log Z(t'), t - t'  \rangle
    \end{split}
\end{equation}
\end{theorem}

\begin{proof}
We begin by applying Bayes' theorem to derive the posterior distribution. Then we use properties of exponential families and the law of large numbers to obtain the desired limit.

By Bayes' theorem:
\begin{equation}
p(t|x_1, \ldots, x_N) = \frac{p(x_1, \ldots, x_N|t)p(t)}{\int_S p(x_1, \ldots, x_N|s)p(s)ds}.
\end{equation}
Since the data samples $x_1, \ldots, x_N$ are i.i.d. given $t$, we have:
\begin{equation}
\begin{split}
p(x_1, \ldots, x_N|t) &= \prod_{i=1}^N p(x_i|t) = \prod_{i=1}^N e^{\langle t, f(x_i) \rangle - \log Z(t)} \\ &= e^{\sum_{i=1}^N \langle t, f(x_i) \rangle - N \log Z(t)}.
\end{split}
\end{equation}
Substituting into the posterior, we find:
\begin{equation}
p(t|x_1, \ldots, x_N)^{1/N} = \frac{e^{\langle t, \frac{1}{N} \sum_{i=1}^N f(x_i) \rangle - \log Z(t) + \frac{1}{N} \log p(t)}}{\left(\int_S e^{N\left[\langle s, \frac{1}{N} \sum_{i=1}^N f(x_i) \rangle - \log Z(s) + \frac{1}{N}\log p(s) \right]} ds \right)^{1/N}}.
\end{equation}

Now, consider the limit:
\begin{equation}
\lim_{N \to \infty} \frac{1}{N} \sum_{i=1}^N f(x_i).
\end{equation}
By Kolmogorov's strong law of large numbers, this converges almost surely to the expectation:
\begin{equation}
\frac{1}{N} \sum_{i=1}^N f(x_i) \xrightarrow{\text{a.s.}} \int f(x) p(x|t') dx = \nabla_{t'} \log Z(t'),
\end{equation}
where the last equality holds because $p(x|t)$ is an exponential family distribution.

In other words:
\begin{equation}
p\left(\{\omega \in \Omega : \lim_{N \to \infty} \frac{1}{N} \sum_{i=1}^N f(x_i(\omega)) = \nabla_{t'} \log Z(t')\} \middle| t' \right) = 1.
\end{equation}

By the continuous mapping theorem: if $X_1, X_2, \ldots$ is a sequence of random variables and $h: \mathbb{R} \to \mathbb{R}$ is a continuous function, then:
\begin{equation}
X_n \xrightarrow{\text{a.s.}} X \implies h(X_n) \xrightarrow{\text{a.s.}} h(X).
\end{equation}
Let $\psi(s, x) = e^{\langle s, x \rangle - \log Z(s) + \log p(s)}$. Since the function $\psi$ is continuous in $x$, and we have shown that $\frac{1}{N} \sum_{i=1}^N f(x_i) \xrightarrow{\text{a.s.}} \nabla_{t'} \log Z(t')$ as $N \to \infty$, the continuous mapping theorem implies:
\begin{equation}
 e^{\langle t, \frac{1}{N} \sum_{i=1}^N f(x_i) \rangle - \log Z(t) + \frac{1}{N}\log p(t)} \xrightarrow{\text{a.s.}} e^{\langle t, \nabla_{t'} \log Z(t') \rangle - \log Z(t)},
\end{equation}
where $p(t)$ is a continuous function on a compact domain, and therefore bounded, which implies $\frac{1}{N} \log p(t) \to 0$ as $N \to \infty$.

Now, let's consider the term in the denominator of the posterior :
\begin{equation}
\left( \int_S e^{N\left[\langle s, \frac{1}{N} \sum_{i=1}^N f(x_i) \rangle - \log Z(s) + \frac{1}{N} \log p(s)\right]} ds \right)^{1/N}.
\end{equation}

We can rewrite the integral part using Lemma \ref{Lemma:Integral}. This lemma applies because, according to the strong law of large numbers, the average $\frac{1}{N} \sum_{i=1}^N f(x_i)$ converges almost surely to a finite value. To utilize the lemma, we define the function:
\begin{equation}
\phi(s, N, \omega) = \langle s, \frac{1}{N} \sum_{i=1}^N f(x_i(\omega)) \rangle - \log Z(s) = \psi\left(s, \frac{1}{N} \sum_{i=1}^N f(x_i(\omega))\right).
\end{equation}
We will ignore the term $\frac{1}{N} \log p(s)$ as it converges to 0 as N goes to infinity and wont affect the value of the integral.
Then by Lemma \ref{Lemma:Integral}, we can state that the following limit holds almost surely:
\begin{equation}
\lim_{N \to \infty} \frac{1}{N} \log \int_S e^{N \phi(s, N, \omega)} ds \overset{\text{a.s.}}{=} \max_{s \in S} \psi(s, \nabla_{t'} \log Z(t')).
\end{equation}
% where we used $\mu_x = \nabla_{t'} \log Z(t')$ from the strong law of large numbers.


Using this result, we have
\begin{equation}
p(t|x_1, \ldots, x_N)^{1/N} = \frac{e^{\langle t, \frac{1}{N} \sum_{i=1}^N f(x_i) \rangle - \log Z(t) + \frac{1}{N} \log p(t)}}{\left(\int_S e^{\langle s, \frac{1}{N} \sum_{i=1}^N f(x_i) \rangle - \log Z(s) + \frac{1}{N} \log p(s)} ds \right)^{1/N}} \xrightarrow{\text{a.s.}} \frac{e^{\langle t, \nabla_{t'} \log Z(t') \rangle - \log Z(t)}}{e^{\max_{s \in S} \psi(s, \nabla_{t'} \log Z(t'))}}.
\end{equation}
Now, observe that we have:
\begin{equation}
\max_{s \in S} \psi(s, \nabla_{t'} \log Z(t')) = \max_{s \in S} \left[ \langle s, \nabla_{t'} \log Z(t') \rangle - \log Z(s) \right].
\end{equation}
Since the function is concave, we can state that the maximum is achieved when $s=t'$, which allows us to obtain
\begin{equation}
\max_{s \in S} \psi(s, \nabla_{t'} \log Z(t')) = \langle t', \nabla_{t'} \log Z(t') \rangle - \log Z(t')
\end{equation}
Substituting back into the limit, we arrive at the expression for the Bregman divergence:
\begin{equation}
\lim_{N \to \infty} p(t|x_1, \ldots, x_N)^{1/N} \overset{\text{a.s.}}{=} \frac{e^{\langle t, \nabla_{t'} \log Z(t') \rangle - \log Z(t)}}{e^{\langle t', \nabla_{t'} \log Z(t') \rangle - \log Z(t')}} = e^{-D_{B}(t',t)}.
\end{equation}
To complete the proof, we need to show that the KL divergence $D_\text{KL}(p(x|t') \| p(x|t))$ in the theorem statement is equivalent to the Bregman divergence $D_B(t, t')$ in our result. We compute the KL divergence as follows:
\begin{align*}
    D_\text{KL}(p(x|t') \| p(x|t)) &= \int p(x|t') \log\frac{p(x|t')}{p(x|t)} dx \\
    &= \int p(x|t') [\log p(x|t') - \log p(x|t)] dx \\
    &= \int p(x|t') [\langle t', f(x) \rangle - \log Z(t') - (\langle t, f(x) \rangle - \log Z(t))] dx \\
    &= \int p(x|t') \langle t', f(x) \rangle dx - \log Z(t') - \int p(x|t') \langle t, f(x) \rangle dx + \log Z(t) \\
    &= \langle t', \nabla_{t'} \log Z(t') \rangle - \log Z(t') - \langle t, \nabla_{t'} \log Z(t') \rangle + \log Z(t) \\
    &= \langle t' - t, \nabla_{t'} \log Z(t') \rangle + \log Z(t) - \log Z(t') \\
    &= D_B(t, t')
\end{align*}
where we use the property that the expectation of $f(x)$ under the distribution $p(x|t')$ is equal to the gradient of the log-partition function $ \mathbb{E}_{x\sim p(x|t')}[f(x)] = \nabla_{t'} \log Z(t')$.

Thus, we have shown that the limit of the posterior is related to the Bregman divergence of $t'$ and $t$, and is consistent with the theorem statement, which concludes the proof.

\end{proof}

\end{document}
